
\section{Methodological Approach and Structure}
The review follows structured IS/medical literature guidance (concept
matrix inspired by \parencite{websterAnalyzingPrepareFuture2002} and PRISMA
principles). The methodological approach uses a multi‑stage search and
transparent extraction/synthesis procedures as outlined below.

\begin{enumerate}
  \item \textbf{Search strategy (Multi-stage, 2015--2025):}
    \begin{itemize}
      \item \textit{Scoping:} Google Scholar for broad coverage and
        identifying candidate papers.
      \item \textit{Formal queries:} IEEE Xplore, PubMed, and Scopus
        using controlled keyword strings and deduplication.
      \item \textit{Citation chasing:} Backward and forward citation
        tracking on key papers to ensure completeness.
    \end{itemize}
  \item \textbf{Inclusion / Exclusion criteria:}
    \begin{itemize}
      \item \textit{Inclusion:} Studies using wearable/non-EEG autonomic
        signals for seizure detection or preictal prediction.
      \item \textit{Exclusion:} EEG-only studies, invasive recordings,
        or papers limited to purely algorithmic simulations without
        human data.
    \end{itemize}
\end{enumerate}


Extracted dimensions populate a concept matrix with the following
fields:
\begin{enumerate}
  \item Setting: clinical vs ambulatory.
  \item Modality: ECG/HRV, PPG, EDA, ACC.
  \item Task: detection vs prediction.
  \item Model class and fusion strategy.
  \item Study phase: retrospective, pseudo‑prospective, prospective.
  \item Evaluation granularity: sample‑ vs alarm‑based.
  \item Windowing specifics (window length, stride, overlap).
  \item Metrics: sensitivity, FAR/h, AUC, f1-score, latency.
  \item Personalization.
  \item Interpretability and safety considerations.
  \item Data characteristics and dataset identity.
  \item Missing data handling (imputation, exclusion) and synthetic data use.
  \item Inference requirements (latency, computational load, hardware).
\end{enumerate}

\noindent\textbf{Prioritization Criteria:} 
Studies including splits preserving patient-specific data, FAR/h reporting 
and external or cross validation will be prioritized in the synthesis. A PRISMA-style flow diagram
will summarize and visualize screening and inclusion.

